% !PRINTED 7
\weberpage{7}

\begin{center}
\S.\,2.

\gesperrt{Die Divisoren endlicher Gruppen.}
\end{center}

Es ist, wie wir schon im ersten Bande gesehen haben, eine besonders wichtige Frage, ob ein Theil der Elemente einer Gruppe selbst wieder eine Gruppe ist. Dazu ist nothwendig und hinreichend, dass je zwei Elemente dieses Theiles bei der Zusammensetzung wieder ein demselben Theile angehöriges Element ergeben. Diese kleinere Gruppe heisst dann ein \gesperrt{Theiler} oder \gesperrt{Divisor}\footnote{Auch Untergruppe genannt.} der ganzen Gruppe. Das Einheitselement „1“ bildet in jeder Gruppe für sich eine Gruppe, ist also in jeder Gruppe ein Divisor.

Wir beschränken uns jetzt, wie in der Folge fast immer, auf \gesperrt{endliche Gruppen}, und wollen annehmen, es sei
\[
P = 1, a_1, a_2 \ldots a_{n-1}
\]
eine Gruppe vom Grade $n$ mit den Elementen $a_0 = 1$, $a_1$, $a_2 \ldots a_{n-1}$ und
\[
Q = 1, a_1, a_2 \ldots a_{\nu-1}
\]
sei ein Theiler von $P$ vom Grade $\nu$. Jeder solche Theiler $Q$ muss die Grundeigenschaften einer Gruppe haben und enthält also sicher das Einheitselement „1“ und mit jedem seiner Elemente $a_1$ zugleich das entgegengesetzte $a_1^{-1}$. Ist $\nu < n$, so nennen wir $Q$ einen \gesperrt{echten Theiler} von $P$. Ist dann also $b$ ein nicht in $Q$ enthaltenes Element von $P$, so sind die Elemente
\[
Q_1 = b, a_1 b, a_2 b \ldots a_{\nu-1} b
\]
alle von einander verschieden (\S.\,1, 3.) und alle nicht in $Q$ enthalten; denn wenn etwa $a_1 b$ in $Q$ enthalten wäre, so müsste auch, da $Q$ eine Gruppe ist, $a_1^{-1} a_1 b = b$ in $Q$ enthalten sein, gegen die Voraussetzung. Ist mit $Q$ und $Q_1$ die ganze Gruppe $P$ noch nicht erschöpft, so nehmen wir eines der noch übrigen Elemente, das wir mit $c$ bezeichnen können, und bilden
\[
Q_2 = c, a_1 c, a_2 c \ldots a_{\nu-1} c,
\]
und überzeugen uns leicht, dass die Elemente von $Q_2$ alle nicht nur von einander, sondern auch von den Elementen von $Q$ und
%CONT>>
